\section{Fibred categories over interval types}
\label{sec:B}

We use $C=\Ty^{\op}$. Its objects are finite words in $+,-$, including
the empty word; its arrows are nondecreasing, sign-preserving maps of
letter sets. Write $P=R^{\op}$ for its reachability preorder and
$p:C\to P$ for the projection. The run number $\nu(w)$ is the number
of maximal constant-sign runs, with $\nu(\emptyset)=0$. An arrow
$w\to v$ implies $\nu(w)\leq\nu(v)$. For equal positive run numbers
it also forces equality of the first signs. Conversely these conditions
are sufficient, with either first sign allowed when the inequality is
strict. Thus reduction of a nonempty word to its alternating run word
describes its class in $P$, without asserting that run reduction is a
functor on $C$.

\begin{claim}{B.1}{Known}{Audits of R12 and the two INDLOCAL comparisons}
\label{B:audits}
R12, INDLOCAL C.21 and INDLOCAL C.23 have the stated variance and
their conclusions used below hold. In particular $C$ is neither
total nor exact, and a projective type sequence is cofinal exactly
when its run numbers are unbounded. The comparison
$E_t\to C_\infty(t)$ is always a weak equivalence. The other leg
is not a weak equivalence uniformly for all diagrams.
\end{claim}
\begin{proof}
For R12, put $a=++$, $b=+-+$ and use the maps
$e,c:a\rightrightarrows b$ with letter maps $(1,3)$ and $(1,1)$.
Every sign-preserving map out of $b$ is injective on its three
letters. Thus every $ge$ is injective and every $hc$ is constant;
no equality $ge=hc$ is possible. On the free module on $C(a,w)$,
the functional taking value one on injective maps and zero on
constant maps satisfies
$\lambda(t(e-c))=\operatorname{aug}(t)$.
This rules out an affine equalizer over every nonzero ring.

We spell out the compressed homological step of TOTAL\_EXACT\_2
C.40, label \texttt{FT:ty}, whose proof we read in full. Let
$P_x=\Z[C(x,-)]$, let $d:P_b\to P_a$ send $g$ to $ge-gc$,
and put $K=\ker d$, $J=\im d$. Every element of $K(w)$ has
augmentation zero. Hence the map
$\colim K\to\colim P_b=\Z$ is zero: a colimit element is a
finite sum of classes represented at possibly different objects,
and each summand has zero augmentation. Right exactness applied
to $K\to P_b\to J\to0$ gives $\colim J=\Z$. Its map into
$\colim P_a=\Z$ is zero, because every $ge-gc$ has augmentation
zero. The derived sequence for
$0\to J\to P_a\to\Z[C(a,-)/(ge\sim gc)]\to0$
therefore gives $T_1=\Z$. The permutation-module bar is exactly
the nerve chains of the element category. This verifies the
integral $H_1$ assertion and finds no error in the return.

In INDLOCAL C.21 the run numbers along a projective sequence are
nondecreasing. Unbounded alternating strings contain every finite
sign word as a weakly increasing sign-preserving letter map.
Conversely maps from all alternating words force unbounded run
number. For a bounded sequence the eventual run word is fixed,
and the contractible hom-posets of types identify its path fibres
with that fixed run word modulo natural transformations. These are
the precise steps retained here; we do not reopen its path-diagram
classification.

For C.23, with $u:J\to T$ and $F:T^{\op}\to\Cat$, the comma of
the induced total-category functor over $(t,z)$ retracts onto
$u/t$. Over $(j,\alpha:u(j)\to t)$ its fibre is the slice at
$F(\alpha)z$, with terminal arrow $\id$. Selecting those arrows
is a natural section, and the given slice arrows contract the
fibres. Thus contractible $u/t$ suffice by Theorem A; the discrete
representable $T(-,t)$ proves necessity. For a sequence, the nerve
of a filtered categorical colimit is its filtered simplicial
colimit, and a filtered homotopy colimit maps weakly equivalently
to that colimit. This proves the first comparison leg, with its
actual map, as in PATHS A.2. The representable of $+$ gives a
contractible all-type total and a sequential discrete colimit
with several elements, so the second leg fails in general.
\end{proof}

\begin{claim}{B.2}{Proved}{Every bounded reachability interval is contractible}
\label{B:interval}
For any types $r\leq s$ in $R$, the full subcategory of $\Ty$
on $\{\tau:r\leq\tau\leq s\}$ is weakly contractible.
\end{claim}
\begin{proof}
Pass to the opposite category, so the interval is $s\leq w\leq r$
in $P$. If $s$ is empty, the interval has initial object the empty
word. Otherwise let $k=\nu(s)>0$ and $l=\nu(r)\geq k$, and replace
only their reachability labels by the alternating run words.

Let $K$ be the nerve of the indiscrete groupoid on two objects
labelled $+,-$. The nonempty-word part of $C$ is exactly
$\Delta/K$: a word lists the vertices of a simplex, and its
monotone sign-preserving maps are precisely simplex maps. Let
$A_r\subset K$ be the image of the alternating $(l-1)$-simplex
whose word is $r$. Its simplices are exactly the words $w\leq r$.
If $k=l$, let $L$ be its boundary subcomplex, consisting of all
lower-dimensional nondegenerate simplices. If $k<l$, let $L$ be
the image of the alternating $(k-1)$-simplex with first sign
opposite to that of $s$. In either case, the simplices of $A_r$
outside $L$ are exactly those with $s\leq w\leq r$.

We establish below two precise geometric facts: the image of an
alternating $d$-simplex of $K$ realizes as a closed $d$-ball with
all lower simplices in its boundary; and the category of simplices
of a simplicial set $A$ outside a simplicial subset $L$ has the
weak type of $|A|\setminus|L|$. Consequently our interval has the
weak type of a closed ball with a closed subset of its boundary
removed. Straight-line contraction to an interior point contracts
that complement. This also covers $k=l$, where the complement is
the open top cell, and $k=l=1$, where it is a point.

% Supplement prepared by the Part G agent for the parent's B.4 proof.
% No claim number is assigned here; integrate into the numbered B statement.
\paragraph{The alternating simplex is a ball.}
Let $K$ be the nerve of the indiscrete groupoid on the two objects
$+,-$, and let $K_n$ be its $n$-skeleton.  Its nondegenerate
$j$-simplices are the two alternating strings of $j+1$ signs.
We prove that $|K_n|$ is an $n$-sphere, and that the image in $K$
of either alternating $n$-simplex realizes as a closed $n$-ball
whose boundary is $|K_{n-1}|$.  In particular any simplicial subset
of that image which omits its top nondegenerate simplex lies in
the boundary of this ball.

Represent a simplex with sign string $(\epsilon_0,\ldots,\epsilon_j)$
and barycentric coordinates $(\lambda_0,\ldots,\lambda_j)$ by the
step function on $[0,1]$ which has value $\epsilon_i\in\{-1,1\}$
on an interval of length $\lambda_i$, the intervals being placed in
the indicated order.  Deleting a zero coordinate and merging
consecutive equal signs are precisely the realization identifications.
It follows that $|K_n|$ identifies bijectively with the almost
everywhere equivalence classes of $\{-1,1\}$-valued step functions
having at most $n$ sign changes in $(0,1)$.  The map to this space
with its $L^1$ metric is continuous on each simplex: moving the
breakpoints changes the function only on intervals whose total
length tends to zero.  Since $K_n$ has finitely many nondegenerate
simplices, its realization is compact.  The $L^1$ target is Hausdorff,
so this bijection is a homeomorphism.

We give a finite-dimensional convex proof of the asserted sphere and
ball descriptions.  For $d\geq1$, put
\[
 p_c(t)=\sum_{j=0}^{d-1}c_jt^j,
 \qquad h_d(c)=\int_0^1|p_c(t)|\,dt,
 \qquad
 M_d=\{x\in\mathbb R^d:c\cdot x\leq h_d(c)
                    \text{ for every }c\in\mathbb R^d\}.
\]
This is a compact convex body with $0$ in its interior.  Indeed the
coordinate inequalities bound it, and the positive continuous function
$h_d$ has a positive minimum on the Euclidean unit sphere, so a small
ball about $0$ lies in $M_d$.  For nonzero $c$, define
\[
 x(c)_j=\int_0^1t^j\operatorname{sign}(p_c(t))\,dt
       \quad(0\leq j<d).
\]
The polynomial has finitely many zeros.  For every $a$ we have
$a\cdot x(c)\leq h_d(a)$, while
$c\cdot x(c)=h_d(c)$.  Thus $h_d$ is exactly the support function
of $M_d$, and $x(c)$ belongs to its support face with normal $c$.

That support face has only one point.  Dominated convergence gives,
for every vector $a$,
\[
 \left.\frac{d}{ds}\right|_{s=0}h_d(c+sa)
   =\int_0^1p_a(t)\operatorname{sign}(p_c(t))\,dt
   =a\cdot x(c).
\]
If $y$ belongs to the same support face, the inequalities
$h_d(c+sa)\geq(c+sa)\cdot y$ for positive and negative $s$
force $a\cdot y=a\cdot x(c)$ for every $a$, hence $y=x(c)$.
Every boundary point has a supporting hyperplane, so every boundary
point of $M_d$ is the moment vector of the sign of a polynomial of
degree at most $d-1$.

Take $d=n+1$.  Any step function $f$ with at most $n$ sign changes
is, almost everywhere, the sign of a nonzero polynomial of degree
at most $n$: use the product of the linear factors at its actual
change points and choose the sign of the leading coefficient.
Consequently its moment vector
\[
 \mu_n(f)=\left(\int_0^1t^j f(t)\,dt\right)_{0\leq j\leq n}
\]
lies on $\partial M_{n+1}$, and the preceding support argument proves
surjectivity onto that boundary.  This map is also injective.
If $f=\operatorname{sign}(p)$ and $\mu_n(g)=\mu_n(f)$, then
\[
 \int_0^1p(t)g(t)\,dt=\int_0^1|p(t)|\,dt.
\]
The integrand $|p|-pg$ is nonnegative, so equality implies $g=f$
almost everywhere.  Continuity and compactness therefore identify
$|K_n|$ homeomorphically with $\partial M_{n+1}$, which is an
$n$-sphere by radial projection from $0$.

It remains to identify the closure of either top simplex, rather
than merely the whole skeleton.  For $n\geq1$, project $M_{n+1}$
onto its first $n$ coordinates.  Its image is $M_n$: their support
functions agree, because $h_{n+1}(c,0)=h_n(c)$.  Write its fibers
as intervals
\[
 \{x\}\times[\ell(x),u(x)],\qquad x\in M_n.
\]
For $x\in\operatorname{int}(M_n)$ the interval has positive length.
For example, since $0\in\operatorname{int}(M_{n+1})$, express $x$
as a strict convex combination of $0$ and another point of $M_n$;
combining a lift of the latter point with a small ball about $0$
gives a nontrivial vertical interval over $x$.  A support normal
to $(x,u(x))$ must then have positive last coefficient: a negative
coefficient is contradicted by the lower point in the same fiber,
and a zero coefficient would support $M_n$ at its interior point
$x$.  Its polynomial therefore has degree $n$ and positive leading
coefficient.  Its sign function must have exactly $n$ sign changes,
for otherwise the first $n$ moments already lie on $\partial M_n$
by the result just proved for $n-1$.  Such a sign function ends
with $+$; conversely every step function with exactly $n$ changes
and ending with $+$ is exposed by a degree-$n$ polynomial with
positive leading coefficient and therefore belongs to this upper
endpoint graph.  The lower endpoint graph similarly consists in
its interior of the functions with exactly $n$ changes ending
with $-$.

For $x\in\partial M_n$, the whole fiber is a singleton: a support
normal $c$ of $M_n$ at $x$ exposes a single point of $M_{n+1}$
with normal $(c,0)$, and every point of the fiber lies in that
support face.  The endpoint functions $u,\ell$ are continuous.
They are respectively concave and convex and hence continuous on
the interior of the convex body.  At a boundary point, compactness
of $M_{n+1}$ and the singleton-fiber assertion force every convergent
subsequence of endpoint values to have the unique boundary value;
boundedness then gives continuity there as well.  Thus each endpoint
graph is homeomorphic by projection to the closed ball $M_n$, and
their common boundary is the image of $|K_{n-1}|$.

Finally the image of either alternating $n$-simplex consists of its
open top simplex together with all of $K_{n-1}$: its two end faces
are the two alternating $(n-1)$-simplices, and all lower simplices
occur recursively.  We have therefore identified exactly that image
with one of the two endpoint graphs.  For $n=0$ the image is a
single point, which supplies the initial case.

\paragraph{Consequence for a removed lower simplicial subset.}
If $A$ is the image of an alternating $n$-simplex and $L\subset A$
is a simplicial subset omitting its top nondegenerate simplex, then
$|L|$ is a closed subset of the boundary of this closed ball.
After a homeomorphism with a convex closed ball, straight-line
contraction to an interior point stays outside $|L|$ at every
positive time.  Hence $|A|\setminus|L|$ is contractible, including
the case $L=\partial A$.  This proves the geometric ingredient of
B.4 without requiring $L$ itself to be a ball.

\paragraph{The relative simplex-category comparison.}
For any simplicial set $A$ and simplicial subset $L\subseteq A$,
let $\mathcal C=\Delta/A$ and let $\mathcal C_{\mathrm{out}}$ be
its full subcategory on the simplices which do not belong to $L$.
There is a specified zigzag of weak equivalences
\[
 B\mathcal C_{\mathrm{out}}
 \ \longleftarrow\ \hocolim_{\mathcal C_{\mathrm{out}}}U
 \ \longrightarrow\ |A|\setminus|L|,
 \qquad
 U(\sigma)=|\Delta^m|\setminus|\sigma^{-1}L|
 \quad(\sigma:\Delta[m]\to A).
\]
We prove this for arbitrary small simplicial sets, including
$A=L$ and $L=\emptyset$.

First extend $U$ to all of $\mathcal C$ by the same formula;
$U(\sigma)$ is empty exactly for $\sigma\in L$.  If
$\alpha:\sigma\to\tau$ is a simplex-category arrow, its affine
map takes $U(\sigma)$ to $U(\tau)$ because
$\tau\alpha=\sigma$.  More precisely the preimage of $|L|$ in
$|\Delta^m|$ is exactly $|\sigma^{-1}L|$.  Indeed a point has a
unique expression in the interior of a nondegenerate simplex, and
it belongs to $|L|$ exactly when that nondegenerate simplex belongs
to $L$.  On a standard simplex this tests the face given by the
positive barycentric coordinates.  Thus the removed set is a union
of closed faces, lying in the boundary whenever $\sigma\notin L$.
Straight-line contraction to the barycenter proves that every
nonempty $U(\sigma)$ is contractible.

The category $\mathcal C$ is Reedy: degree is simplex dimension,
direct arrows are injective operators, inverse arrows are surjective
operators, and the unique image factorization of an order-preserving
map supplies the Reedy factorization.  Its inverse matching category
at a nondegenerate simplex is empty.  At a degenerate simplex it
has terminal object the unique surjection onto its nondegenerate
core: every other degeneracy representation factors uniquely through
this one by the Eilenberg--Zilber decomposition.  Hence $\mathcal C$
has fibrant constants.  These are precisely Hirschhorn's
Example~15.1.19 and Proposition~15.10.4(1), with the proof just
checked; Theorem~15.10.8(2) says that colimit is then left Quillen
for the Reedy model structure \cite{BHirschhorn2003}.

The diagram $U$ is Reedy cofibrant in the usual model category of
compactly generated spaces.  Its latching map at $\sigma$ is
\[
 |\partial\Delta^m|\setminus|\sigma^{-1}L|
    \ \lhook\joinrel\longrightarrow\
 |\Delta^m|\setminus|\sigma^{-1}L|.
\]
To verify an actual cofibration, rather than merely a homotopy-type
assertion, put $Z=|\sigma^{-1}L|$.  If $Z$ is empty we have the
usual finite CW pair, and if $Z$ is the whole simplex both spaces
are empty.  Otherwise use barycentric coordinates and the continuous
piecewise linear function
\[
 \delta(x)=
 \max_{\substack{\emptyset\ne S\subseteq\{0,\ldots,m\}\\
                  \Delta^S\not\subseteq Z}}
       \min_{i\in S}x_i.
\]
It vanishes exactly on $Z$: a support face belongs to $Z$ if and
only if every nonempty subset of that support belongs to $Z$.
The dyadic shells
$2^{-k-1}\leq\delta\leq2^{-k}$, together with the uppermost
shell, are finite polyhedra.  Cutting by all simplex faces and
these level sets and taking compatible barycentric subdivisions
triangulates the complement locally finitely, because every point
of the complement has $\delta>0$.  In this triangulation its
boundary complement is a subcomplex.  Thus the displayed map is
a relative CW inclusion.  The finite union of proper faces is
exactly the latching object, with its subspace topology; this also
follows by gluing its finite closed face cover.  Hence all latching
maps are cofibrations, as asserted.

Hirschhorn's Theorem~19.9.1(1), applied with the checked hypotheses,
now gives a weak equivalence
\[
 \hocolim_{\mathcal C}U\longrightarrow\colim_{\mathcal C}U.
\]
The colimit on the right is precisely $|A|\setminus|L|$.
The usual realization quotient has $U(\sigma)$ as the preimage
of this open complement on each simplex, and restricting a quotient
map to the inverse image of an open subset is a quotient map.
This proves both the set and topology assertions, with no finiteness
restriction on $A$.

Finally $\Delta/L\subseteq\mathcal C$ is a sieve.  In the bar
formula for this homotopy colimit, a nonempty summand is indexed
by a chain $\sigma_0\to\cdots\to\sigma_q$ with
$U(\sigma_0)\ne\emptyset$.  Such a chain cannot enter the sieve,
so it lies entirely in $\mathcal C_{\mathrm{out}}$.  Hence the
homotopy colimit over $\mathcal C$ is exactly that over
$\mathcal C_{\mathrm{out}}$.  The objectwise contraction maps
$U\to *$ induce a weak equivalence to
$\hocolim_{\mathcal C_{\mathrm{out}}}*=B\mathcal C_{\mathrm{out}}$
by homotopy invariance, Hirschhorn's Theorem~19.4.2(1).
This supplies the announced zigzag.  For $A=L$ all three spaces
are empty, so that edge case is included.

\end{proof}

\begin{claim}{B.3}{Proved}{The three local coherence conditions are equivalent}
\label{B:coherence}
For every strict $F:C\to\Cat$, with the conventions of the
question,
\[
                 (2F)\quad\Longleftrightarrow\quad
                 (\mathrm{DESC})\quad\Longleftrightarrow\quad
                 (\mathrm{LOC}).
\]
The equivalence includes the allowed zigzags in $(2F)$; it does
not require a chosen strict action on the given representative $F$.
\end{claim}
\begin{proof}
We check the ingredient retained from PATHS C.P1. For two types,
the poset of endpoint-fixed monotone interval functors is empty or
contractible. In the nonempty case remove an initial source run
of the wrong sign; every type map is constant there. For matching
initial signs, move every noninitial image vertex up to the last
vertex of the first target run when necessary. All moved vertices
lie on that same run, so this is an endofunctor with a natural
transformation to or from the identity according to the sign.
Its image has the first run prescribed. Delete the two prescribed
runs and induct. This proves the claimed contraction with its
directions checked. Composition makes these posets the strict
2-category $\TTy$.

A strict action on $\TTy^{\op}$ therefore gives, after applying
classifying spaces, an enriched diagram whose indexing
hom-spaces are empty or contractible. The projection of this
enriched category to $P$ is the identity on objects and an
equivalence on every hom-space; it is a Dwyer--Kan equivalence.
Restriction along it is an equivalence on homotopy-coherent
space-valued diagrams. Thus $(2F)$ implies $(\mathrm{DESC})$.
Objectwise weak-equivalence zigzags preserve this conclusion.
This is the precise coherent descent in PATHS C.P2, rather than
an inference from equality of maps in an ordinary homotopy category.

For $r\in C$, put $C_{\leq r}=\{w:p(w)\leq p(r)\}$.
The homotopy left Kan extension of $BF$ along $p$ has value
\[
       (p_!BF)(r)=\hocolim_{C_{\leq r}}BF.
\]
The inclusion of the fibre in $(\mathrm{LOC})$ is its unit, under
the natural Grothendieck-construction comparison. If $BF=p^*G$
coherently, the functor $C_{\leq r}\to P_{\leq r}$ is homotopy
final: its comma at $s\leq r$ is precisely the interval of B.2.
Hence this homotopy colimit is $G(r)$, since $r$ is terminal in
$P_{\leq r}$. The unit is a weak equivalence. This proves
$(\mathrm{DESC})\Rightarrow(\mathrm{LOC})$.

For the converse we supply an explicit strictification. Use
covariant Grothendieck constructions and define strict diagrams
\[
 G(r)=\int_{C_{\leq r}}^{\rm co}F\quad(r\in P),\qquad
 H(r)=\int_{C/r}^{\rm co}F\circ\mathrm{source}\quad(r\in C).
\]
Maps of $G$ are inclusions of full indexing subcategories;
maps of $H$ postcompose the displayed arrow into $r$.
There are strict natural functors
\[
 F\ \xleftarrow{q}\ H\ \xrightarrow{k}\ p^*G.
\]
At $r$, $q$ sends $(w,\alpha:w\to r,x)$ to $F(\alpha)x$;
$k$ forgets $\alpha$. The section $x\mapsto(r,\id,x)$ is
right adjoint to $q$, as the arrow $(\alpha,\id)$ supplies
the adjunction unit. Thus $q_r$ and its section are weak
equivalences. The composite of that section with $k_r$ is the
fibre inclusion into $G(r)$, a weak equivalence by $(\mathrm{LOC})$.
Consequently $k_r$ is a weak equivalence as well. The covariant
and fibred versions used here have the same natural homotopy
colimit model; applying opposites to the fibres and to the total
construction checks the variance. Finally $p^*G$ extends strictly
to $\TTy^{\op}$ by sending every 2-cell to an identity. The
displayed zigzag proves $(2F)$ exactly in its stipulated sense.
\end{proof}

\subsection{An intrinsic criterion without quantifying over sequences}

Let $\Omega$ be the category of countable linearly ordered sets
labelled by $+,-$ which contain alternating finite subsets of
arbitrarily large size. Its maps are all nondecreasing
sign-preserving functions. We work with a set of concrete countable
models closed under finite subsets and tagged ordered sums; this
makes all indexing categories small. Fix
$J=(\mathbb N,+,-,+,-,\ldots)\in\Omega$.
For a countable labelled order $L$ define the category
\[
 F^\omega(L)=\colim_{S\subset L\text{ finite}}F(S),
 \tag{B.4.1}
\]
where finite $S$ is given its induced order and signs and identified
with its unique standard word. For a monotone sign-preserving map
$f:L\to L'$, send the $S$-stage through the actual finite map
$S\to f(S)$ and then into the colimit. This specifies a functor,
including maps which identify equal-sign letters.

\begin{claim}{B.4}{Proved}{Complete classification of the comparison property}
\label{B:classification}
A strict diagram $F:\Ty^{\op}\to\Cat$ satisfies $(\mathcal D)$
if and only if $F^\omega$ takes every map in $\Omega$ to a weak
equivalence. Equivalently, for every $L\in\Omega$, both specified
inclusions in
\[
 F^\omega(L)\longrightarrow F^\omega(L+J)
                 \longleftarrow F^\omega(J)
 \tag{B.4.2}
\]
are weak equivalences. Here $+$ is ordered concatenation, and
the arrows are induced by its two canonical inclusions.
Thus (B.4.2) is a necessary and sufficient condition on $F$ alone,
with no choice of transition maps or cofinal sequence in its statement.
\end{claim}
\begin{proof}
First, every projective type sequence, viewed covariantly in $C$,
has a filtered colimit $L$ in countable labelled linear orders.
Its order is well defined on the equivalence classes: a finite
list of representatives can be compared at a common stage, and
distinct limiting classes remain distinct. Signs are unchanged
by all transition maps. Each alternating finite subset of a stage
remains distinct in the limit, so cofinality is equivalent to
$L\in\Omega$. Conversely enumerate any $L\in\Omega$ and take
the increasing sequence of its finite subsets. Their induced
words form a cofinal sequence with colimit $L$.

Every finite map into a filtered colimit of such orders factors
at a stage, and an equality of two finite maps holds at a later
stage. This follows by putting their finitely many representatives
and finitely many equality conditions at one common stage.
Consequently the left Kan extension (B.4.1) gives a canonical
isomorphism of categories
\[
       \colim_n F(t_n)\ \cong\ F^\omega(L).
\]
The assertion remains valid for noninjective transitions.

Let $\widehat C$ be the category of all finite or countable labelled
linear orders in our chosen set of models. The category of finite
subsets of $L$ is homotopy final in $C/L$: over a finite map
$a\to L$, its comma is the filtered poset of finite subsets
containing its image, with that image as initial object. Hence
$BF^\omega$ is also the homotopy left Kan extension of $BF$
from $C$ to $\widehat C$. Filtered categorical colimits and their
homotopy colimits agree here by B.1.

The inclusion $\Omega\to\widehat C$ is homotopy final. In the
comma over $L$, apply $T(A,f)=(A+J,\iota_A f)$. There are natural
arrows from the identity and from the constant object
$(L+J,\iota_L)$ to $T$, the latter induced by $f+\id_J$.
These contract the comma. The same sum operation, with the
inclusions of $L$ and $J$, contracts $B\Omega$ itself. Fubini
for homotopy left Kan extensions and finality now give a natural
weak equivalence
\[
 \hocolim_C BF\ \simeq\ \hocolim_{L\in\Omega}BF^\omega(L).
 \tag{B.4.3}
\]
Its comparison at a particular $L$ is the canonical roof of
the question, using any finite-subset sequence for $L$. For an
arbitrary sequence with limit $L$, the finite-factorization
argument above identifies its roof with the same comparison.
This identification is natural in every map of labelled orders.

If $(\mathcal D)$ holds, these comparisons are all equivalences;
naturality and two out of three make every $F^\omega(f)$ an
equivalence. Conversely, a diagram of spaces sending every arrow
of $\Omega$ to an equivalence factors through its groupoid
completion. Since $B\Omega$ is contractible, each value maps
equivalently to its homotopy colimit. Equation (B.4.3) proves
$(\mathcal D)$, including its actual canonical comparison.

The all-arrows condition clearly implies (B.4.2). Conversely,
the functor $T(L)=L+J$ has natural transformations
$\id_\Omega\to T\leftarrow\mathrm{const}_J$.
If both components in (B.4.2) are weak equivalences, naturality
first makes $F^\omega(Tf)$ and then $F^\omega(f)$ a weak
equivalence for every $f$. This proves the claimed equivalent
test and both directions of the classification.
\end{proof}

\begin{claim}{B.5}{Proved}{The precise sequential form of Lead 5(b)}
\label{B:infinity}
The map $BE_t\to B\int_{\Ty}F$ is, naturally in the homotopy
category of arrows, the sequential homotopy colimit of the units
\[
 BF(t_n)\longrightarrow B\int_{t_n\uparrow}F.
\]
In particular $(\mathrm{LOC})\Rightarrow(\mathcal D)$.
Deleting any fixed bounded set of run numbers does not change
the homotopy-colimit comparisons for any $F$.
\end{claim}
\begin{proof}
The full subcategories $t_n\uparrow$ of $\Ty$ increase with $n$
and their union is $\Ty$. Their fibred total categories increase
as full subcategories and exhaust $\int F$, including all arrows
and every finite nerve simplex. Thus their sequential homotopy
colimit maps weakly equivalently to the total category. The unit
maps form the natural homotopy-coherent unit of the homotopy left
Kan extension in B.3. Their coherence is also the explicit
cartesian arrow for a transition of types. Taking the homotopy
colimit of this unit gives precisely the comparison, by the
natural Grothendieck-construction theorem. This is an equivalence
of specified arrows in spaces; it is not an assertion that the
two displayed diagrams of categories commute strictly before
those cartesian transformations are included.

For the last assertion let $C_{>m}$ be the full subcategory of
words with more than $m$ runs. For $a\in C$ choose a fixed
alternating buffer $b$ with more than $m+1$ runs. In
$a\backslash C_{>m}$, the endofunctor $(w,f)\mapsto(w+b,\iota f)$
receives natural arrows from the identity and from the constant
$(a+b,\iota_a)$, using $f+\id_b$. All words involved remain
in $C_{>m}$. The comma is contractible, so that full inclusion
is homotopy final for every diagram. Every cofinal sequence also
eventually lies there. This proves the asserted invariance at
infinity without imposing a condition on low fibres.
\end{proof}

\begin{claim}{B.6}{Refuted}{The comparison property need not imply local coherence}
\label{B:DnotLOC}
There is an explicit presheaf of sets satisfying $(\mathcal D)$
but failing all three conditions in B.3.
\end{claim}
\begin{proof}
Put $a=++$ and define the covariant set diagram on $C$
\[
 S(w)=
 \begin{cases}C(a,w),&\nu(w)\leq1,\\ \{*\},&\nu(w)>1.
 \end{cases}
\]
On the lower region use postcomposition, on the upper region use
the identity of a singleton, and across the boundary use the
unique map. There are no arrows from the upper region to the
lower, so this defines all maps and satisfies composition.
By B.5 its total homotopy colimit is the contractible nerve of
$C_{>1}$; that category is contracted by $w\mapsto w+b$ with
the two block inclusions. Every cofinal sequence is eventually
the singleton diagram, so $(\mathcal D)$ holds. At $a$, however,
$S(a)$ has three elements. The downset $C_{\leq a}$ consists
of the empty word and the nonempty positive words; its element
category for $C(a,-)$ has initial object $(a,\id_a)$.
The local comparison is therefore a three-point space mapping
to a contractible space and fails on components. B.3 gives the
other two failures. The empty and all-negative lower fibres
are empty exactly as the representable formula specifies.
\end{proof}

\begin{claim}{B.7}{Refuted}{One cofinal sequence does not suffice}
\label{B:D1notD}
There is a finitely presented set diagram $S:C\to\Set$ with
contractible all-type total category such that its comparison
is a weak equivalence along a dense cofinal sequence and fails
on components along the alternating sequence $J$.
Consequently $(\mathcal D_1)$ does not imply $(\mathcal D)$,
even for presheaves of sets.
\end{claim}
\begin{proof}
Take $a=+-$ and $b=+-+-$, and let $i:a\to b$ select the first
and last letters. Put $R=C(a,-)$ and let $T\subset R$ be the
image of $C(b,-)\to C(a,-)$ induced by $i$. Thus an element of
$R(w)$ is a marked pair $p<q$ with signs $+,-$; it belongs to
$T(w)$ precisely when there are $p<x<y<q$ with signs $-,+$.
All four images are distinct by sign preservation and monotonicity.
Set
\[
                 S=R\amalg_T1_C.
 \tag{B.7.1}
\]
This is the pointed quotient collapsing all pairs in $T(w)$
to the added basepoint. It is also the pushout of
$C(b,-)\to C(a,-)$ and $C(b,-)\to C(\emptyset,-)$, so is
finitely presented. These formulas specify every transition map.

We prove $B\El_C T$ contractible. First retain only its
injective underlying arrows, obtaining $E_{\rm inj}$. Its
objects are finite signed orders with marked endpoints $p,q$
and at least one extra $-,+$ pair between them. Amalgamate a
variable such object with the fixed marked word $b$, identifying
the two marked endpoints. In each of the three gaps cut out
by those endpoints, put variable-side points before fixed-side
points. On an injective arrow fixing the endpoints use that
arrow on the variable side and the identity on $b$. No outside
point can enter a marked endpoint under such an injection.
This gives an endofunctor with natural arrows from the identity
and from the constant $b$. Thus $BE_{\rm inj}$ is contractible.

The inclusion $j:E_{\rm inj}\to\El_C T$ is a weak equivalence.
Fix a target marked word $w$. An object of $j/w$ is a monotone
sign-preserving map $g:v\to w$ preserving the marked endpoints;
its arrows are injections over $w$. Amalgamate $v$ with a fixed
copy of $w$ over those endpoints and map the union to $w$ by
$g$ and the identity. Order the union first by its image letter
in $w$. In a nonmarked image fibre put all variable-side points
before the fixed point. In either marked fibre put the variable
points preceding its marked point first, then the identified
marked point, then the remaining variable points. This order
extends both source orders and preserves signs. An injection
over $w$ preserves each image fibre and each side of a marked
point, so the construction is functorial on $j/w$. It receives
natural injections from the variable object and from $(w,\id)$.
It stays in $T$ because it contains the fixed $w$. Hence every
such comma is contractible, and Theorem A proves the assertion.

The objectwise pushout square (B.7.1) is a homotopy pushout of
spaces, since $T(w)\to R(w)$ is injective. Homotopy colimits
preserve this square. The three homotopy colimits for $T,R,1_C$
are contractible: the first was just proved, the second is the
representable element category with initial object $(a,\id)$,
and the third is $BC$, with initial empty word. Consequently
$\hocolim_C S$ is contractible.

For the good sequence take the dyadic rationals with their usual
order. Give a reduced dyadic rational sign $+$ or $-$ according
as its denominator exponent is even or odd; integers have sign
$+$. Each sign class is dense. Enumerate the set and take the
induced finite-subset sequence. Its run numbers are unbounded.
Between every $+$ point $p$ and $-$ point $q>p$ lie a $-$ point
and then a $+$ point. Hence every element of $R(L)$ belongs
to $T(L)$, and $S^\omega(L)=\{*\}$. The comparison between
this singleton and the contractible total is a weak equivalence.

For the bad sequence use the initial finite subsets of $J$.
Its limit has
\[
 S^\omega(J)=\{*\}\amalg
       \{(2m,2m+1):m\in\mathbb N\}.
\]
Indeed a $+$--$-$ pair at distance at least three has the
intermediate $-,+$ witness, while an adjacent pair never does.
Thus the comparison fails on components. The program
\texttt{code/B\_type\_checks.py} independently checks the
finite presentation, functor laws, the two amalgam ordering
rules, and the adjacent-pair invariant. The proofs above, rather
than a finite bound, establish the two infinite-sequence claims.
\end{proof}

\begin{claim}{B.8}{Proved}{All implications and the exact answer for sets}
\label{B:table}
For arbitrary category-valued diagrams, and also within
presheaves of sets, the full implication diagram is
\[
 (2F)\ \Longleftrightarrow\ (\mathrm{DESC})
 \ \Longleftrightarrow\ (\mathrm{LOC})
 \ \Longrightarrow\ (\mathcal D)
 \ \Longrightarrow\ (\mathcal D_1),
\]
and neither final implication reverses. Every implication
between two of the five conditions follows exactly by
transitivity from this display.

For a set diagram $S$, the first three conditions are equivalent
to $S$ factoring through $P$, up to its unique transported natural
isomorphism. The diagrams for which $\Ty/S$ has weakly
contractible components indexed by $\colim_n S(t_n)$ for every
cofinal sequence are exactly those for which the two set maps
in (B.4.2) are bijective for every $L\in\Omega$.
\end{claim}
\begin{proof}
B.3 and B.5 prove the positive implications. B.6 refutes each
reverse implication from $(\mathcal D)$ to a condition in the
first equivalence class. B.7 refutes $(\mathcal D_1)$ implying
any stronger condition. Both examples are set diagrams, so
the same failures hold in that subcategory.

For discrete values, a natural transformation along a 2-cell
forces equality of its two maps. Connectedness of every
nonempty type hom-poset therefore forces all parallel maps
to have the same effect. This is precisely factorization
through $P$. Alternatively take components of a coherent
descent equivalence; a coherent diagram of discrete spaces is
an ordinary set diagram, and its equivalences are bijections.
Conversely such a factorization gives a strict action with
all 2-cells identities.

For the last assertion apply B.4 to discrete categories. The
filtered categorical colimit is discrete, so a weak equivalence
is exactly a bijection. If its criterion holds, the extension
to $\Omega$ is a constant diagram of sets up to a natural
isomorphism, since $B\Omega$ is contractible and all transition
maps are bijective. Equation (B.4.3) identifies the total with
that discrete set, canonically on components. Conversely the
specified component and contractibility assertion says exactly
that every comparison from its discrete sequential colimit is
a weak equivalence, so B.4 gives the criterion. This proof
includes empty presheaves, where all the spaces are empty.
\end{proof}

\begin{claim}{B.9}{Proved}{The comparison as strict 2-categorical weak equivalences}
\label{B:two-category}
Let $F:\TTy^{\op}\to\Cat$ be a strict 2-functor. Regard
ordinary categories as locally discrete strict 2-categories,
and write $\mathcal E_F$ for its fibred 2-Grothendieck construction.
For every cofinal sequence $t$, both arrows in the specified roof
\[
       C_\infty(t)\ \longleftarrow E_t
                         \longrightarrow\mathcal E_F
\]
are weak equivalences in the Thomason structure used in 317.
The comparison $\int_{\Ty}F\to\mathcal E_F$, identity on
objects and 1-arrows, is likewise a strict 2-functor and a weak
equivalence. These statements are invariant under the
objectwise weak-equivalence replacements allowed in $(2F)$.
\end{claim}
\begin{proof}
The 2-construction has objects $(\tau,x)$, 1-arrows
$(\phi,f):(\tau,x)\to(\sigma,y)$ with
$f:x\to F(\phi)y$, and a 2-arrow induced by
$\alpha:\phi\Rightarrow\psi$ precisely when
$F(\alpha)_y f=g$. Composition and the interchange law are
those of the strict 2-action. Thus the two stated functors
are defined on all cells.

The enriched homotopy-colimit theorem
\cite[Theorem 5.1 in the opened preprint]{BCegarraHeredia2015}
identifies its classifying
space with the enriched bar of $BF$. Under the equivalence of
enriched indexing categories with $P$ used in B.3, this bar is
the homotopy colimit of its descended diagram on $P$. The
ordinary bar is the restriction along $C\to P$. This functor
is homotopy final: over a fixed type, concatenating it with a
variable predecessor gives a contraction by the two block
projections in $\Ty$, or their inclusions in $C$. The sequence
is homotopy final over $P$ because its comma at each type is
a nonempty final segment of $\mathbb N$. Therefore the actual
ordinary-to-enriched comparison and the sequence comparison
are weak equivalences. Naturality of the bar construction
identifies these as the displayed functors, not abstract
isomorphisms of their homotopy types. The other roof leg is B.1.

For the precise 2-categorical convention we checked 317,
Section 7, equation \texttt{eq:model-classes}: its weak
equivalences are exactly those detected by the Street nerve.
On ordinary categories that nerve is the ordinary nerve, as
verified directly in its foundations section. In dimension two,
the model structure and this class are also the established
Thomason structure of Ara--Maltsiniotis,
\cite[Theorem 6.27]{BAraMaltsiniotis2014}.
The enriched-bar/Grothendieck comparison used here is
Cegarra--Heredia's homotopy-colimit theorem; its comparison of
the double nerve and the 2-categorical nerve gives the same
weak equivalences. Only this dimension-two specialization of
317 is needed, so no claim about a supplied Lean certificate
or about higher-dimensional transfer is being used without
verification. Objectwise equivalences induce equivalences
of both bars, proving the final invariance assertion.
\end{proof}

\begin{claim}{B.10}{Proved}{The answer to VII ms. 99--115}
\label{B:manuscript}
At VII ms. 103(b), the sequential model computes the all-type
fibred homotopy type for every cofinal sequence exactly for the
diagrams of B.4. For coherent type actions, the answer is always
positive; that sufficient class is strictly smaller than the
complete class. A positive answer for one cofinal sequence does
not imply a positive answer for every one.
\end{claim}
\begin{proof}
The equivalence is B.4, the sufficient coherent class is B.3,
and strictness and dependence on the sequence are B.6 and B.7.
The canonical roof, rather than an unspecified functor from the
strict colimit to the total category, is the isomorphism in
$\Hot$ asserted here. The representable of $+$ gives the
failure of a uniform all-diagram assertion; the one-object-group
example at ms. 109 is already correctly rejected by the comma
criterion in INDLOCAL C.23. At ms. 109--110 and 112--114,
cofinality of the reachability preorder alone therefore remains
insufficient. The phrase \emph{en g\'en\'eral pas} at ms. 112
is correct, while the full path diagrams are positive because
their strict 2-action supplies B.3. Their one-sign and bounded-run
exceptions are exactly those already classified in INDLOCAL
C.21, as audited in B.1, and are not being claimed cofinal here.
Finally B.2 supplies the interval contractibility omitted from
Lead 5(c), and B.5 supplies the precise sequential homotopy-colimit
meaning of Lead 5(b). This verifies Lead 5(a)--(c), while
refuting the stronger suggestion that $(\mathcal D)$ holds
only for diagrams admitting the coherent type action.
\end{proof}
